\chapter{ALJABAR PENGALI}\label{multiplier_algebras}

\section{Modul Hilbert}
\begin{notn}\label{0038014} Hasil kali dalam yang muncul sebelumnya dalam catatan ini
dan yang dijumpai dalam buku teks dan monograf standar tentang ruang Hilbert,
analisis fungsional, dan sebagainya, bersifat linear dalam variabel pertama dan linear konjugat dalam
variabel kedua. Kebanyakan ahli aljabar operator kontemporer memilih bekerja dengan objek
yang disebut modul Hilbert-$A$ \emph{kanan} ($A$ suatu aljabar-$C^*$). Untuk modul semacam itu,
ternyata lebih mudah menggunakan ``hasil kali dalam'' yang linear dalam variabel kedua
dan linear konjugat dalam variabel pertama. Walaupun perubahan konvensi ini mungkin menimbulkan
sedikit rasa jengkel, secara matematis pengaruhnya kecil. Tentu saja,
kita ingin ruang Hilbert menjadi contoh modul Hilbert-$\C$. Agar hal ini
mungkin, kita membekali ruang Hilbert, yang hasil kali dalamnya dinyatakan dengan
$\langle\hphantom{X},\hphantom{X}\rangle$, dengan ``hasil kali dalam'' baru yang didefinisikan
 \index{<bracketspointy@$\langle x\,\vert\,y \rangle$ (hasil kali dalam)}%
oleh $\langle x\,\vert\,y\rangle := \langle y,x \rangle$. ``Hasil kali dalam'' ini linear
dalam variabel kedua dan linear konjugat dalam variabel pertama. Saya akan berusaha konsisten dalam
menggunakan $\langle\hphantom{X},\hphantom{X}\rangle$ untuk hasil kali dalam standar dan
$\langle\hphantom{X}\,\vert\,\hphantom{X}\rangle$ untuk hasil kali dalam yang linear dalam
variabel keduanya.

Solusi lain (yang sangat umum) untuk persoalan ini adalah menetapkan bahwa hasil kali dalam
selalu linear dalam variabel keduanya dan ``mengoreksi'' buku teks, monograf, dan
makalah standar sesuai dengan ketetapan tersebut.
\end{notn}

\begin{conv} Berdasarkan uraian sebelumnya, selanjutnya kita akan menggunakan kata
``seskuilinear'' untuk berarti salah satu di antara
 \index{seskuilinear}%
 \index{konvensi!tentang \emph{seskuilinear}}%
\emph{linear dalam variabel pertama dan linear konjugat dalam variabel kedua} atau \emph{linear dalam
variabel kedua dan linear konjugat dalam variabel pertama}.
\end{conv}

\begin{defn}\label{0038017} Misalkan $A$ suatu aljabar-$C^*$ tak nol. Suatu ruang vektor $V$ adalah
\df{modul-$A$}
 \index{modul}%
 \index{amodule@modul-$A$}%
jika terdapat suatu pemetaan bilinear
  \[ B\colon V \times A \sto V\colon (x,a) \mapsto xa \]
sedemikian sehingga $x(ab) = (xa)b$ berlaku untuk semua $x \in V$ dan $a$, $b \in A$. Kita juga mensyaratkan
bahwa $x \vc 1_A = x$ untuk setiap $x \in V$ jika $A$ beridentitas. (Suatu fungsi dua variabel
 \index{bilinear}%
disebut \df{bilinear} jika fungsi itu linear dalam kedua variabelnya.)
\end{defn}

\begin{defn}\label{0038021} Untuk memperjelas sebagian materi berikutnya, akan
berguna untuk memiliki definisi formal berikut. Suatu
 \index{vektor!ruang}%
\df{ruang vektor (kompleks)} adalah tripel $(V,+,M)$ dengan $(V,+)$ suatu grup Abel dan
$M\colon \C \sto \Hom(V,+)$ suatu homomorfisme gelanggang beridentitas. (Seperti dapat diduga,
$\Hom(V,+)$ menyatakan gelanggang beridentitas dari endomorfisme grup $(V,+)$.) Jadi, suatu
 \index{aljabar}%
\df{aljabar} adalah kuadruplet terurut $(A,+,M,\,\cdot\,)$ dengan $(A,+,M)$ suatu ruang
vektor dan $\,\cdot\,\colon A \times A \sto A \colon (a,b) \mapsto a \cdot b = ab$ suatu
operasi biner pada $A$ yang memenuhi syarat-syarat dalam~\ref{defn_alg}.
\end{defn}

\begin{exer} Pastikan bahwa definisi \emph{ruang vektor} sebelumnya ekuivalen
dengan definisi yang biasa Anda gunakan.
\end{exer}

\begin{defn} Misalkan $(A,+,M,\,\cdot\,)$ suatu aljabar (kompleks). Maka $A^{\textrm{op}}$
 \index{<A@$A^{\textrm{op}}$ (aljabar lawan dari $A$)}%
adalah aljabar $(A,+,M,*)$ dengan $a*b = b\cdot a$ untuk semua $a$, $b \in A$. Aljabar ini adalah
 \index{aljabar lawan}%
 \index{aljabar!lawan}%
\df{aljabar lawan} dari~$A$. Aljabar itu hanyalah $A$ dengan urutan perkalian dibalik.
Jika $A$ dan $B$ adalah aljabar, maka setiap fungsi $f \colon A \sto B$ dengan cara yang jelas menginduksi
suatu fungsi dari $A$ ke $B^{\textrm{op}}$ (atau dari $A^{\textrm{op}}$ ke $B$,
atau dari $A^{\textrm{op}}$ ke $B^{\textrm{op}}$). Semua fungsi ini akan kita nyatakan
cukup dengan~$f$.
\end{defn}

\begin{defn}\label{0038027} Misalkan $A$ dan $B$ aljabar. Suatu fungsi
$\phi\colon A \sto B$ adalah suatu
 \index{antihomomorfisme}%
\df{antihomomorfisme} jika fungsi $\phi\colon A \sto B^{\textrm{op}}$ merupakan suatu homomorfisme.
Suatu antihomomorfisme bijektif adalah
 \index{antiisomorfisme}%
\df{antiisomorfisme}. Dalam contoh~\ref{0022057} kita menjumpai \emph{antiisomorfisme}
ruang Hilbert. Antiisomorfisme tersebut agak berbeda karena alih-alih membalik perkalian
(yang tidak dimiliki ruang Hilbert), antiisomorfisme itu mengonjugatkan skalar. Dalam kedua kasus,
antiisomorfisme bukanlah sesuatu yang sama sekali berbeda dari isomorfisme, melainkan justru
sesuatu yang sangat serupa.
\end{defn}

Gagasan modul-$A$ (dengan $A$ suatu aljabar) telah didefinisikan dalam~\ref{0038017}. Anda
mungkin lebih menyukai definisi alternatif berikut, yang lebih selaras dengan
definisi \emph{ruang vektor} dalam~\ref{0038021}.

\begin{defn}\label{0038031} Misalkan $A$ suatu aljabar. Kuadruplet terurut
 \index{modul}%
 \index{amodule@modul-$A$}%
$(V,+,M,\Phi)$ disebut \df{modul-$A$} jika $(V,+,M)$ suatu ruang vektor dan $\Phi\colon A \sto \ofml L(V)$ suatu
homomorfisme aljabar. Jika $A$ beridentitas, kita juga mensyaratkan bahwa $\Phi$ beridentitas.
\end{defn}

\begin{exer} Periksa bahwa definisi \emph{modul-$A$} dalam~\ref{0038017}
dan~\ref{0038031} ekuivalen.
\end{exer}

Sekarang kita nyatakan secara tepat apa artinya, ketika $A$ suatu aljabar-$C^*$, membekali modul-$A$ dengan
hasil kali dalam bernilai-$A$.

\begin{defn} Misalkan $A$ suatu aljabar-$C^*$. Suatu
 \index{hasil kali dalam semu!modul-$A$}%
 \index{modul!hasil kali dalam semu modul-$A$}%
 \index{amodule@modul-$A$!hasil kali dalam semu}%
\df{modul-$A$ hasil kali dalam semu} adalah suatu modul-$A$ $V$ beserta pemetaan
  \[ \beta\colon V \times V \sto A\colon (x,y) \mapsto \langle x\vert\, y\rangle \]
yang linear dalam variabel keduanya dan memenuhi
 \begin{enumerate}
    \item[(i)] $\langle x\vert\, ya\rangle = \langle x\vert\, y\rangle a$,
    \item[(ii)] $\langle x\vert\, y\rangle = \langle y\vert\, x\rangle^*$, dan
    \item[(iii)] $\langle x\vert\, x\rangle \ge \vc 0$
 \end{enumerate}
untuk semua $x$, $y \in V$ dan $a \in A$. Modul ini adalah suatu
 \index{hasil kali dalam!modul-$A$}%
 \index{modul!hasil kali dalam modul-$A$}%
 \index{amodule@modul-$A$!hasil kali dalam}%
\df{modul-$A$ hasil kali dalam} (atau suatu
 \index{pra-Hilbert!modul-$A$}%
 \index{modul!pra-Hilbert-$A$}%
 \index{amodule@modul-$A$!pra-Hilbert}%
\df{modul pra-Hilbert-$A$}) jika sebagai tambahan
 \begin{enumerate}
    \item[(iv)] $\langle x\vert\, x\rangle = \vc 0$ mengakibatkan $x = \vc 0$
 \end{enumerate}
untuk $x \in V$. Kita akan menyebut pemetaan $\beta$ sebagai
 \index{hasil kali dalam!bernilai-$A$}%
 \index{Aval@hasil kali dalam (semu) bernilai-$A$}%
\df{hasil kali dalam (semu) bernilai-$A$} pada~$V$.
\end{defn}

\begin{exam} Setiap ruang hasil kali dalam adalah modul-$\C$ hasil kali dalam.
\end{exam}

\begin{prop} Misalkan $A$ suatu aljabar-$C^*$ dan $V$ suatu modul-$A$ hasil kali dalam semu.
Hasil kali dalam semu $\langle\hphantom{x}\vert\,\hphantom{x}\rangle$ bersifat linear konjugat
dalam variabel pertamanya, baik secara harfiah maupun dalam arti bahwa $\langle va \vert\,
w\rangle = a^* \langle v \vert\, w\rangle$ untuk semua $v$, $w \in V$ dan $a \in A$.
\end{prop}

\begin{prop}[Ketaksamaan Schwarz---untuk modul-$A$ hasil kali dalam]\label{0038038} Misalkan
$V$ suatu modul-$A$ hasil kali dalam dengan $A$ suatu aljabar-$C^*$. Maka
  \[ \langle x\vert\, y \rangle^*\langle x\vert\, y \rangle
          \le \norm{\langle x\vert\, x \rangle}\,\langle y\vert\, y \rangle \]
untuk semua $x$, $y \in V$.
\end{prop}

\begin{proof}[\emph{Petunjuk pembuktian}] Tunjukkan bahwa, tanpa mengurangi keumuman,
kita dapat mengasumsikan bahwa
$\norm{\langle x\vert\, x\rangle} = 1$. Pertimbangkan elemen positif $\langle xa -
y\vert\, xa - y\rangle$ dengan $a = \langle x \vert\, y \rangle$. Gunakan
proposisi~\ref{0018203} dan~\ref{0018204}. \ns
\end{proof}

\begin{defn}\label{0038041} Untuk setiap elemen $v$ dari suatu modul-$A$ hasil kali dalam
(dengan $A$ suatu aljabar-$C^*$), definisikan
  \[ \norm v := \norm{\langle v \vert\, v \rangle}^{1/2}. \]
\end{defn}

\begin{prop}[Ketaksamaan Schwarz lainnya] Misalkan $A$ suatu aljabar-$C^*$ dan $V$ suatu
modul-$A$ hasil kali dalam. Maka untuk semua $v$, $w \in V$
  \[ \norm{\langle v \vert\, w\rangle} \le \norm v\,\norm w. \]
\end{prop}

\begin{cor} Jika $v$ dan $w$ merupakan elemen suatu modul-$A$ hasil kali dalam (dengan $A$ suatu
aljabar-$C^*$), maka $\norm{v + w} \le \norm v + \norm w$ dan pemetaan $x \mapsto \norm x$
merupakan norma pada~$V$.
\end{cor}

\begin{prop} Jika $A$ suatu aljabar-$C^*$ dan $V$ suatu modul-$A$ hasil kali dalam, maka
  \[ \norm{va} \le \norm v\, \norm a \]
untuk semua $v \in V$ dan $a \in A$.
\end{prop}

\begin{defn} Misalkan $A$ suatu aljabar-$C^*$ dan $V$ suatu modul-$A$ hasil kali dalam. Jika $V$
lengkap terhadap (metrik yang diinduksi oleh) norma yang didefinisikan dalam~\ref{0038041},
maka $V$ adalah suatu
 \index{amodule@modul-$A$!Hilbert}%
 \index{Hilbert!modul-$A$}%
\df{modul Hilbert-$A$}.
\end{defn}

\begin{exam}\label{0038111} Untuk $a$ dan $b$ dalam suatu aljabar-$C^*$ $A$,
 \index{Hilbert!modul-$A$!$A$ sebagai}%
 \index{amodule@modul-$A$!Hilbert!$A$ sebagai}%
definisikan
  \[ \langle a \vert\, b\rangle := a^*b\,. \]
Maka $A$ sendiri merupakan modul Hilbert-$A$. Setiap ideal kanan tertutup dalam $A$ juga merupakan
modul Hilbert-$A$.
\end{exam}

\begin{defn} Misalkan $V$ dan $W$ modul-modul Hilbert-$A$ dengan $A$ suatu aljabar-$C^*$. Suatu
pemetaan $T\colon V \sto W$ bersifat
 \index{alinear@linear-$A$}%
 \index{linear!$A$-}%
\df{linear-$A$} jika pemetaan itu linear dan jika $T(va) = T(v)a$ berlaku untuk semua $v \in V$ dan $a \in
A$. Pemetaan $T$ adalah suatu
 \index{amodule@modul-$A$!morfisme}%
 \index{modul!morfisme}%
 \index{morfisme!modul Hilbert-$A$}%
 \index{Hilbert!modul-$A$!morfisme}%
\df{morfisme modul Hilbert-$A$} jika pemetaan itu terbatas dan linear-$A$.
\end{defn}

Ingat dari~\ref{def000926} bahwa setiap operator ruang Hilbert memiliki adjoin.
Hal ini tidak berlaku untuk modul Hilbert-$A$.

\begin{defn} Misalkan $V$ dan $W$ modul-modul Hilbert-$A$ dengan $A$ suatu aljabar-$C^*$. Suatu
fungsi $T\colon V \sto W$ bersifat
 \index{dapat diadjoinkan}%
\df{dapat diadjoinkan} jika terdapat suatu fungsi $T^*\colon W \sto V$ yang memenuhi
  \[ \langle Tv \vert\, w \rangle = \langle v \vert\, T^*w \rangle \]
untuk semua $v \in V$ dan $w \in W$. Fungsi $T^*$, jika ada, adalah
 \index{adjoin}
\df{adjoin} dari~$T$. Nyatakan dengan
 \index{L@$\ofml L(V,W)$, $\ofml L(V)$ (pemetaan yang dapat diadjoinkan)}%
$\ofml L(V,W)$ keluarga pemetaan yang dapat diadjoinkan dari $V$ ke~$W$. Kita menyingkat $\ofml L(V,V)$
menjadi~$\ofml L(V)$.
\end{defn}

\begin{prop} Misalkan $V$ dan $W$ modul-modul Hilbert-$A$ dengan $A$ suatu aljabar-$C^*$. Jika suatu
fungsi $T\colon V \sto W$ dapat diadjoinkan, maka fungsi itu merupakan morfisme modul Hilbert-$A$.
Lebih lanjut, jika $T$ dapat diadjoinkan, maka adjoinnya juga dapat diadjoinkan dan $T^{**} = T$.
\end{prop}

\begin{exam} Misalkan $X$ interval satuan $[0,1]$ dan $Y = \{0\}$. Dengan topologi biasa,
$X$ merupakan ruang Hausdorff kompak dan $Y$ suatu subruang dari~$X$. Misalkan $A$ adalah
aljabar-$C^*$ $\fml C(X)$ dan $J_0$ adalah ideal $\{f \in A\colon f(0) = 0\}$ (lihat
proposisi~\ref{0006836}). Pandang $V = A$ dan $W = J_0$ sebagai modul-modul Hilbert-$A$ (lihat
contoh~\ref{0038111}). Maka pemetaan inklusi
$\iota\colon W \sto V$ merupakan morfisme modul Hilbert-$A$ yang tidak dapat diadjoinkan.
\end{exam}

\begin{prop} Misalkan $A$ suatu aljabar-$C^*$. Pasangan pemetaan $V \mapsto V$, $T \mapsto T^*$
merupakan funktor kontravarian dari kategori modul Hilbert-$A$ dan pemetaan yang dapat diadjoinkan
ke dirinya sendiri.
\end{prop}

\begin{prop}\label{0038244} Misalkan $A$ suatu aljabar-$C^*$ dan $V$ suatu modul Hilbert-$A$.
 \index{C*@$C^*$-aljabar!$\ofml L(V)$ sebagai suatu}%
 \index{L@$\ofml L(V)$!sebagai suatu aljabar-$C^*$}%
Maka $\ofml L(V)$ merupakan aljabar-$C^*$ beridentitas.
\end{prop}

\begin{notn} Misalkan $V$ dan $W$ modul-modul Hilbert-$A$ dengan $A$ suatu aljabar-$C^*$.
 \index{theta@$\Theta_{v,w}$}%
Untuk $v \in V$ dan $w \in W$, misalkan
  \[ \Theta_{v,w}\colon W \sto V\colon x \mapsto v\langle w \vert\, x\rangle\,. \]
(Bandingkan dengan~\ref{0022431fa}.)
\end{notn}

\begin{prop} Pemetaan $\Theta$ yang didefinisikan di atas bersifat seskuilinear.
\end{prop}

\begin{prop} Misalkan $V$ dan $W$ modul-modul Hilbert-$A$ dengan $A$ suatu aljabar-$C^*$. Untuk
setiap $v \in V$ dan $w \in W$, pemetaan $\Theta_{v,w}$ dapat diadjoinkan dan $(\Theta_{v,w})^*
= \Theta_{w,v}$.
\end{prop}

Proposisi berikut memperumum proposisi~\ref{prop_op_act_otimes1} dan~\ref{prop_op_act_otimes2}.

\begin{prop} Misalkan $U$, $V$, $W$, dan $Z$ modul-modul Hilbert-$A$ dengan $A$ suatu aljabar-$C^*$.
Jika $S \in \ofml L(Z,W)$ dan $T \in \ofml L(V,U)$, maka
  \[ T \Theta_{v,w} = \Theta_{Tv,w} \qquad \text{ dan } \qquad
                               \Theta_{v,w} S = \Theta_{v,S^*w} \]
untuk semua $v \in V$ dan $w \in W$.
  \[ Z \to^S W \to^{\Theta_{v,w}} V \to^T U \]
\end{prop}

\begin{prop} Misalkan $U$, $V$, dan $W$ modul-modul Hilbert-$A$ dengan $A$ suatu aljabar-$C^*$.
Andaikan $u \in U$; $v$, $v' \in V$; dan $w \in W$. Maka
  \[ \Theta_{u,v}\Theta_{v',w} = \Theta_{u\langle v \vert\,v' \rangle, w}
                               = \Theta_{u, w\langle v'\vert\,v \rangle}\,. \]
\end{prop}

\begin{notn}\label{0038331} Misalkan $A$ suatu aljabar-$C^*$ dan $V$ serta $W$
modul-modul Hilbert-$A$. Kita nyatakan dengan
 \index{K@$\ofml K(V,W)$, $\ofml K(V)$}%
$\ofml K(W,V)$ rentang linear tertutup dari $\{\Theta_{v,w}\colon v \in V \text{ dan } w \in
W\}$. Seperti biasa, kita menyingkat $\ofml K(V,V)$ menjadi~$\ofml K(V)$.
\end{notn}

\begin{prop}\label{0038334} Jika $A$ suatu aljabar-$C^*$ dan $V$
suatu modul Hilbert-$A$, maka $\ofml K(V)$ adalah ideal dalam aljabar-$C^*$~$\ofml L(V)$.
\end{prop}

Contoh berikut dimaksudkan sebagai pembenaran bagi praktik standar untuk mengidentifikasi suatu
aljabar-$C^*$ $A$ dengan $\ofml K(A)$.

\begin{exam}\label{0038337} Jika suatu aljabar-$C^*$ $A$ kita pandang sebagai modul-$A$ (lihat
contoh~\ref{0038111}), maka $\cstariso{\ofml K(A)}A$.
\end{exam}

\begin{proof}[\emph{Petunjuk pembuktian}] Seperti dalam akibat~\ref{0013151}, untuk setiap
$a \in A$ definisikan operator perkalian kiri $L_a\colon A \sto A\colon x \mapsto ax$. Tunjukkan
bahwa setiap operator semacam itu dapat diadjoinkan dan bahwa pemetaan $L\colon A \sto \ofml L(A) \colon
a \mapsto L_a$ merupakan isomorfisme-$*\,$ ke suatu subaljabar-$C^*$ dari~$\ofml L(A)$. Kemudian
verifikasikan bahwa $\ofml K(A)$ adalah tutupan citra oleh $L$ dari rentang linear hasil kali
elemen-elemen~$A$. \ns
\end{proof}

\begin{exam}\label{0038341} Misalkan $H$ suatu ruang Hilbert yang dipandang sebagai modul-$\C$. Maka
$\ofml K(H)$ (sebagaimana didefinisikan dalam~\ref{0038331}) adalah ideal operator kompak pada~$H$ (lihat
proposisi~\ref{prop_K_is_clo_FR}).
\end{exam}

\begin{proof} Lihat~\cite{RaeburnW:1998}, contoh 2.27. \ns \end{proof}

Contoh sebelumnya telah mendorong banyak peneliti, ketika menangani suatu modul Hilbert sebarang~$V$,
menyebut elemen $\ofml K(V)$ sebagai \emph{operator kompak}. Istilah ini
meragukan karena operator semacam itu belum tentu kompak. (Sebagai contoh, jika
suatu aljabar-$C^*$ beridentitas berdimensi tak hingga $A$ kita pandang sebagai modul-$A$, maka
$\Theta_{\vc 1,\vc 1} = I_A$, tetapi operator identitas pada $A$ tidak kompak---lihat
contoh~\ref{cor_id_op_not_cpt}.)

Fakta bahwa modul Hilbert-$C^*$ hanya digunakan secara terbatas dalam catatan ini hendaknya
tidak membuat Anda berpikir bahwa kajiannya sangat khusus dan/atau kurang penting.
Sebaliknya,
kajian tersebut sekarang merupakan bidang penelitian penting dan dinamis yang memiliki penerapan
pada bidang-bidang yang beragam seperti teori-$K$, aljabar-$C^*$ graf, grup kuantum, probabilitas
kuantum, berkas vektor, geometri nonkomutatif, topologi aljabar dan topologi geometri,
ruang dan aljabar operator, serta wavelet. Lihatlah
laman web Michael Frank~\cite{Frank:2010}, \emph{Hilbert C*-modules and related subjects---a guided
reference overview}, tempat ia mencantumkan 1531 referensi (pada pembaruan 11.09.10) berupa buku,
makalah, dan tesis tentang modul semacam itu dan mengelompokkannya menurut penerapan. Terdapat
grafik menarik (pada halaman 9) yang menggambarkan pertumbuhan bidang
matematika ini. Grafik tersebut mencakup materi dari upaya perintis pada tahun 1950-an dan awal 1960-an
(0--2 makalah per tahun) hingga saat catatan ini ditulis (sekitar 100 makalah per tahun).





























